\begin{abstract}
\looseness=-1
Compiling GUI procedures that agents execute repeatedly into programs can reduce their token costs.
However, measuring payback and deciding when to compile have two challenges.
First, compilation costs are uncertain because attempts can require repair and still fail to produce a usable program.
Second, future reuse is unknown because tasks may stop arriving or GUI drift may stop the program from working.
To address these challenges, we propose PACE (Payback-Aware Compilation from Experience), a system with a measurement protocol and an online compilation algorithm.
The measurement protocol records successful and failed compilation costs, and compares agent and program execution costs on matched task inputs to estimate per-use savings and payback counts.
Using these measurements, the online algorithm compares estimated future savings with compilation costs, including failed attempts, based on past task arrivals and compilation outcomes.
It checks execution and compilation charges against a cumulative budget determined by observed task arrivals before allowing either action.
Under stated action-cost assumptions, total cost after each arrival is at most $1+\epsilon$ times the cost of running every task with the agent.
For successful compilation attempts, estimated payback counts excluding source agent runs are $2$--$16$ uses.
In simulations using recorded task arrivals, PACE reduces token costs by $17.3\%$ compared with ReAct, $24.9\%$ with the AutoRPA adaptation, and $17.3\%$ with the ToolPro adaptation on average ($\epsilon=0.25$).
\end{abstract}
